\documentclass[11pt,a4paper]{article}
\usepackage[T1]{fontenc}
\usepackage{lmodern,amsmath,amssymb,mathtools,booktabs,microtype}
\usepackage[margin=22mm]{geometry}
\usepackage{xcolor,fancyhdr}
\usepackage[colorlinks=true,linkcolor=blue!45!black,urlcolor=blue!45!black]{hyperref}
\hypersetup{pdftitle={Exact maps and analytic certificates: key advances},pdfauthor={KokunoYumeto; AI-assisted research record},pdfsubject={Results-first companion, 6 September 2026}}
\setlength{\parindent}{0pt}
\setlength{\parskip}{5pt}
\setlength{\headheight}{14pt}
\pagestyle{fancy}\fancyhf{}
\fancyhead[L]{\small Exact maps and analytic certificates}
\fancyhead[R]{\small 6 September 2026}
\fancyfoot[L]{\small Advances companion; full proofs in the project archive}
\fancyfoot[R]{\thepage}
\newcommand{\Z}{\mathbb Z}\newcommand{\C}{\mathbb C}
\newcommand{\R}{\mathbb R}\newcommand{\OO}{\mathcal O}
\newcommand{\VV}{\mathcal V_F}\newcommand{\Le}{\Lambda_{\mathrm L}}
\newcommand{\proofloc}[1]{\par\smallskip{\small\textit{Full proof:} #1}\par}
\begin{document}
\thispagestyle{plain}
{\LARGE\bfseries Exact maps and analytic certificates\par}
{\large Key advances in the continuing reconstruction\par}
6 September 2026\quad\textbar\quad AI-assisted mathematical research record

The starting point is the claimed complex-threefold construction on $S^6$
circulated by Levent Alp\"oge with Fable. The continuing audit has produced
explicit results about its analytic fibration and about the marked
Niemeier--Leech--triality constructions investigated alongside it.
This companion gives the actual formulas and proof locations for a selection
of those advances. It does not replace the complete arguments.

\section*{1. The anticanonical ring recovers the fibration}
Let $f:X\to B=\mathbb P^1$ denote the analytic construction with its original
periods and fillings. Its finite reduced fibres $S_1,S_2$ and base coordinate
$t$ satisfy
\[
 f^*(p_1)=3S_1,\qquad f^*(p_2)=4S_2,\qquad
 t(p_1)=0,\quad t(p_2)=1,\quad t(p_0)=\infty.
\]
The evaluation $f^*f_*\omega_{X/B}\to\omega_{X/B}$ has zero divisor
\emph{exactly} $2S_2$. The resulting identities are
\[
 \omega_X\simeq\OO_X(-2S_2),\qquad
 \omega_X^{-2}\simeq f^*\OO_B(p_2),
\]
\[
 f_*\OO_X(kS_2)=\OO_B\bigl(\lfloor k/4\rfloor p_2\bigr)
 \qquad(k\in\Z).
\]
The last equality is equality of subsheaves of meromorphic functions:
on the original cover $t_2=s_2^4$, an invariant Laurent series contains
only powers divisible by four. The pole condition is precisely
$4\operatorname{ord}_{p_2}(b)+k\geq0$.

Consequently the \emph{full graded algebra}, not merely its dimensions, is
\[
 \bigoplus_{m\geq0}H^0(X,\omega_X^{-m})\simeq\C[U,V],
 \qquad \deg U=1,\quad\deg V=2.
\]
Here $U\mapsto u$ with $\operatorname{div}(u)=2S_2$, and
$V\mapsto v=f^*\ell$ under the specified degree-two identification,
where $\ell\in H^0(B,\OO_B(p_2))$ is nonzero at $p_2$.
In degree $m$, the basis is
$u^{m-2j}v^j$ for $0\leq j\leq\lfloor m/2\rfloor$.
Writing $\ell=A/(t-1)+B'$ with $A\ne0$, the complete degree-two map is
\[
 x\longmapsto [t(f(x))-1:A+B'(t(f(x))-1)].
\]
Thus the fibration is recovered intrinsically from $X$, up to this explicitly
retained projective-coordinate choice. This is what makes it possible to test
a deformation of the whole threefold, rather than just a marked torus fibre.
\proofloc{\emph{Complete current workbench}, Theorem 53.17, pp.~699--700;
source \texttt{analytic\_canonical\_ring.tex}, equations CR1--CR5.}

\newpage
\section*{2. The original period constant gives a nonzero deformation}
Keep the cusp-regular period solution $\beta^{\rm part}$ and put
\[
 \beta_c=\beta^{\rm part}+c,\qquad
 D_{\rm part}=\operatorname{Im}\beta^{\rm part}
       -\frac{6(\operatorname{Im}\mu)^2}{\operatorname{Im}\tau},
 \qquad M=\max_{B\setminus\{p_0\}}D_{\rm part}.
\]
On every open disc $Q$ whose closure lies in
$\{c\in\C:\operatorname{Im}c<-M\}$, the original charts and gluing maps
give a proper holomorphic submersion $\mathcal X_Q\to Q$ with fibres $X_c$.
The construction retains all three critical values $0,1,\infty$.

At a fixed lifted regular base point, fix $c_*\in Q$ and write
\[
 T=\operatorname{Im}\tau>0,\quad m=\operatorname{Im}\mu,\quad
 D=\operatorname{Im}\beta_{c_*}-6m^2/T<0,\quad
 P=\begin{pmatrix}0&0\\-m/T&1\end{pmatrix},\quad\delta=c-c_*.
\]
The unchanged ordered period matrix is
\[
 \Pi_c=\begin{pmatrix}6\mu&\tau&1&0\\\beta_c&\mu&0&1\end{pmatrix},
 \qquad \Pi_c:\R^4\longrightarrow\C^2.
\]
Its inverse below is real-linear. The exact marked map between the two period lattices is
\[
 w=\left(I+\frac{\delta}{2iD}P\right)\zeta
            -\frac{\delta}{2iD}P\overline\zeta,
 \qquad \det_{\R}(\Pi_c\Pi_{c_*}^{-1})
       =\frac{D+\operatorname{Im}\delta}{D}>0.
\]
In the convention that the transported $(0,1)$ fields are
$\partial_{\bar\zeta_k}+\sum_j\nu^j_k\partial_{\zeta_j}$,
\[
 \nu_\delta=\frac{\delta}{2iD+\delta}P.
\]
With the connecting-class convention $\mathrm{KS}(\partial_c)=[\bar\partial h]$,
the regular-fibre Kodaira--Spencer class is
\[
 -\left[\frac{1}{2iD}\partial_{\zeta_2}\otimes
       \left(d\bar\zeta_2-\frac mT d\bar\zeta_1\right)\right]\ne0.
\]
Its constant nonzero coefficient cannot be a derivative of a periodic
coefficient: integrating such a derivative over the torus gives zero.

The additional \emph{global} step is essential. Over
$A=\C[\epsilon]/(\epsilon^2)$, the intrinsic degree-two anticanonical
sections form a free $A$-module of rank two. A first-order trivialization of
$X_c$ reducing to the identity at $c_*$ would identify the fibrations through an element of
$\mathrm{PGL}_2(A)$. The actual critical schemes fix the three constant
sections $0,1,\infty$, forcing that element to be the identity. Restriction
to a regular fibre then contradicts the displayed nonzero class. Hence
\[
 \mathrm{KS}_{X_{c_*}}(\partial_c)\ne0\quad\hbox{in }H^1(X_{c_*},T_{X_{c_*}}).
\]
This does not assert that every two parameters give nonisomorphic unmarked
manifolds, or that the family exhausts all deformations.
\proofloc{Workbench, Theorem 53.18, pp.~701--705;
\texttt{period\_parameter\_deformation.tex}, PD1--PD18.}

\newpage
\section*{3. Finite fillings: explicit covers and the full attachment kernel}
At the two original finite centres retain
\[
 \rho=e^{\pi i/3},\quad \mu_1=(2-\rho)/3,\quad\mu_2=(1-i)/2,
 \qquad\kappa_1=\beta_1+2\mu_1,\quad\kappa_2=\beta_2+3\mu_2,
\]
with $\operatorname{Im}\kappa_j<0$. Set
\[
 \begin{gathered}
 E_j=\C/(\Z+\kappa_j\Z),\\
 F_1=\C/(\Z(\rho-2)+\Z(-3)),\qquad
 F_2=\C/(\Z(i-1)+\Z(-2)).
 \end{gathered}
\]
The maps from $E_j\times F_j$ to the original central torus are induced by
\[
 L_1(\xi,\eta)=(\eta,\xi-\eta/3),\qquad
 L_2(\xi,\eta)=(\eta,\xi-\eta/2).
\]
Their kernels are generated respectively by $(2/3,\rho)$ and $(1/2,i)$,
of orders three and two. The lifted affine generators are
\[
 g_1(\xi,\eta)=(\xi+\kappa_1/3,(\rho-1)\eta),\qquad
 g_2(\xi,\eta)=(\xi-\kappa_2/4,-i\eta).
\]
Together with the kernel translations they act freely, giving product
covers of the central surfaces of degrees nine and eight. The varying
cover is identified with the normalization of the \emph{actual} ramified
base change, including every branch of $s_j^{m_j}=u^{m_j}$.

Let $\Lambda$ now denote the original rank-four fibre lattice, and let
$A_j,v_j,m_j$ be its finite linear monodromy, translation vector and order,
with $(m_1,m_2)=(3,4)$ and $A_jv_j=v_j$.
The complete deck-group presentation and attachment kernel are
\[
 \Gamma_j=\langle\Lambda,G_j\mid
 G_jt_\lambda G_j^{-1}=t_{A_j\lambda},\ G_j^{m_j}=t_{v_j}\rangle,
\]
\[
 \ker\bigl(\Lambda\rtimes_{A_j}\Z\longrightarrow\Gamma_j\bigr)
   =\langle\widetilde G_j^{m_j}t_{-v_j}\rangle\simeq\Z,
\]
and this kernel is central. The computation determines the entire kernel,
not merely a relation that lies in it.

Finally let $\gamma\in\Lambda^*$ be the retained invariant coordinate,
$\varepsilon_j=\gamma(v_j)\in\{1,-1\}$ and
$\zeta_j=e^{-2\pi i/m_j}$. The normal line $N_{S_j/N_j}$ has exact
holomorphic order $m_j$, but the explicit nonzero smooth section is
\[
 H_j(x)=\exp(-2\pi i\varepsilon_j\gamma(x)),\qquad
 H_j(A_jx+v_j/m_j)=\zeta_jH_j(x).
\]
The smooth trivialization is
$[x,n]\mapsto([x],n/H_j(x))$, with inverse
$([x],z)\mapsto[x,zH_j(x)]$.
\proofloc{Workbench, Theorems 53.4, 53.6, 53.7 and 53.9,
pp.~684--690; \texttt{finite\_filling\_certificates.tex}, FF1--FF46.}

\newpage
\section*{4. Niemeier--Leech transport retains the cubic arithmetic}
Start with the marked root lattice $R=A_5^{\perp4}\perp D_4$ and its
Niemeier overlattice $N$, with $[N:R]=72$. Its rank is $24$;
$72$ is the glue index, not a dimension.
For its original Weyl vector $\rho_N$ and retained root $\alpha$, put
\[
 K=\{x\in N:(\rho_N,x)\equiv0\pmod6\},\qquad
 \Le=K+\Z(\rho_N/6-\alpha).
\]
Then $N\cap\Le=K$ and $[N:K]=[\Le:K]=6$. The residue
order-three action cycles three such neighbours. Their common pairwise
intersection $J$ has determinant $81$ and index $9$ in each neighbour.
The fourth extension $M$, fixed by this action, has root system $A_2^{12}$
and $72$ roots. These are explicit integral constructions, retaining the
original residue marking and glue generators.

Use the fixed real isometry $\mathcal L$ to $\mathbb O^3$ with the original
Cayley frame $(1,i,j,k,\ell,i\ell,j\ell,k\ell)$. Put
\[
 \tau(q)=\operatorname{Re}((q_1q_2)q_3),\qquad
 F=\tau\circ\mathcal L,\qquad
 \VV(\mathcal A)=\langle F(x):x\in\mathcal A\rangle_{\Z}.
\]
For a lattice $\mathcal A$, this defines an \emph{additive group generated by values};
it does not identify the nonlinear image $F(\mathcal A)$ with that group.
For the retained cubic, the complete calculation gives
\[
\begin{aligned}
 \VV(N)&=\tfrac1{27}\Z\oplus\tfrac{\sqrt2}{54}\Z
                         \oplus\tfrac{\sqrt3}{18}\Z,\\
 \VV(J)=\VV(M)&=\tfrac1{27}\Z\oplus\tfrac{\sqrt2}{216}\Z
                         \oplus\tfrac{\sqrt3}{18}\Z,\\
 \VV(\Le)&=\tfrac1{27}\Z\oplus\tfrac{\sqrt2}{216}\Z
                         \oplus\tfrac{\sqrt3}{54}\Z.
\end{aligned}
\]
In the displayed ordered generators the two group inclusions have matrices
\[
 \VV(N)\hookrightarrow\VV(J):\ \operatorname{diag}(1,4,1),\qquad
 \VV(J)\hookrightarrow\VV(\Le):\ \operatorname{diag}(1,1,3).
\]
Their quotients are $\Z/4\Z$ and $\Z/3\Z$; the total quotient is
$\Z/12\Z$. These are inclusions of \emph{value groups}, not a claim that
$N\subset\Le$. The proof gives both integral containment and reverse
generator certificates using complete cubic polarization.

For the same embedding into the Albert algebra, the equal-diagonal
determinant is
\[
 \det Q_{\mathbb O}(\lambda;\mathcal Ly)
       =\lambda^3-\lambda(y,y)+2F(y).
\]
The coefficient $2$ remains part of the map. Beyond the lattice, with
$E(q)=Q_{\mathbb O}(0;q)$, the spectral graph
\[
 q\longmapsto Q_{\mathbb O}(-\lambda_{\min}(E(q));q)
\]
maps
onto the equal-diagonal positive-semidefinite determinant boundary,
with inverse given by its ordered off-diagonal coordinates. The associated
Peirce products retain their coefficient $1/2$.
\proofloc{\emph{Higher torus wrapping and the octonionic 24-dimensional
continuation}, Sections 9--11, pp.~106--125; spectral graph Theorem 5.12,
p.~34. Source labels: \texttt{nca:sec:cubic-arithmetic},
\texttt{esrf:sec:fixed-neighbor}, \texttt{jmv:sec:JM-extension}.}

\newpage
\section*{5. A literal common lattice in the original triality space}
Wilson's octonionic Leech construction is an established input, not a new
construction claimed here. Retain its norm
$q_W(x)=\tfrac12\sum_{j=1}^3|x_j|^2$ and the isometry
$\iota(x)=x/\sqrt2$ into the Euclidean norm $\sum|q_j|^2$.
Write $\Lambda_{\rm cyc}=\iota(\Lambda_W)$.
The new comparison computes both literal intersections:
\[
 \Lambda_{\rm cyc}\cap\mathcal L(N)
 =\Lambda_{\rm cyc}\cap\mathcal L(\Le)
 =\mathcal I
 =\{\sqrt2(a,a,a):a\in D_4\},
\]
where here $D_4=\{a\in\Z^4:\sum_{i=0}^3a_i\equiv0\pmod2\}$ lies
in the \emph{first four original Cayley coordinates}. The Gram matrix,
determinant and minimum squared norm of $\mathcal I$ are
\[
 6G_{D_4},\qquad 5184,\qquad 12.
\]
The cubic on this common lattice is explicit:
\[
 \tau(\sqrt2(a,a,a))
 =2\sqrt2\,\underbrace{a_0\bigl(a_0^2-3(a_1^2+a_2^2+a_3^2)\bigr)}_{P(a)}.
\]
The generated group is $4\sqrt2\Z$, but $12\sqrt2$ is not attained.
Indeed $P(a)\equiv a_0^3\pmod3$, so $3\mid P(a)$ forces $3\mid a_0$;
substitution then gives $9\mid P(a)$. The required value $P(a)=6$ is impossible.
This gives a concrete arithmetic distinction between a group of values
and individual values, while retaining their exact common embedding.
\proofloc{Higher-rung paper, Section 12, pp.~125--135;
\texttt{es\_niemeier\_wilson\_cyclic.tex}.
Primary input: R.~A.~Wilson, \emph{Octonions and the Leech lattice},
J.~Algebra \textbf{322} (2009), 2186--2190,
\href{https://doi.org/10.1016/j.jalgebra.2009.03.021}{published paper};
\href{https://webspace.maths.qmul.ac.uk/r.a.wilson/pubs_files/octoLeech1rev.pdf}{author version},
Sections 2--4.}

\section*{Reading the release}
The \emph{complete higher-rung paper} contains the triality and lattice
proofs. The \emph{complete current workbench} also contains the analytic
ring, deformation and finite-filling arguments. Printed page numbers above
refer to the 6 September copies (138 and 782 pages respectively).

The separately labelled \emph{complete-project calculations and provenance
ZIP} retains the full TeX, exact checkers and results, mathematical datasets,
standalone papers and mathematical history, including work not selected for
this front. The older published files remain unchanged as dated records.
All proof-source filenames above lie under
\texttt{project/supporting\_materials/workbench/research/} in that archive.

These advances do not by themselves certify every global step in the
$S^6$ claim or establish a counterexample to CDP20. No interacting
quantum-field-theory mass-gap conclusion follows from the displayed lattice
maps. The complete workbench retains the arguments and their dependencies
for further checking; this is an AI-assisted research record, not independent
expert adjudication or Lean formalization.

\href{https://doi.org/10.5281/zenodo.22235523}{Release family and full inventory}
\quad\textbar\quad
\href{https://www.overleaf.com/read/rtmyqxyrzprn#fa24eb}{Continuing readable workspace}
\end{document}
